\chapter{Podciągi, ciągi Cauchy'ego i zupełność}
\label{ma1:ch:cauchy-completeness}

W poprzednim rozdziale badaliśmy zbieżność ciągu do zadanego punktu. Teraz rozpatrzymy dwa zagadnienia: jak wydobyć z ciągu zbieżny podciąg i jak rozpoznać zbieżność bez znajomości granicy. Pierwsze prowadzi do twierdzenia Bolzana--Weierstrassa, drugie do warunku Cauchy'ego oraz pojęcia zupełności. Własności ogólne formułujemy metrycznie, a własności porządkowe osobno dla liczb rzeczywistych. Związki ze zwartością i zupełnością w topologii omawiane są szerzej w Topologii ogólnej.

\section{Podciągi i punkty skupienia ciągu}
\begin{definicja}
\label{ma1:def:subsequence}
Niech $(x_n)_{n\ge1}$ będzie ciągiem w zbiorze $X$. Podciągiem nazywamy ciąg $(x_{n_k})_{k\ge1}$, gdzie $(n_k)_{k\ge1}$ jest ściśle rosnącym ciągiem dodatnich liczb naturalnych. Jeśli $(X,d)$ jest przestrzenią metryczną, punkt $a\in X$ nazywamy punktem skupienia ciągu $(x_n)$, jeżeli pewien jego podciąg jest zbieżny do $a$.
\end{definicja}
Pojęcie to dotyczy ciągu z uwzględnieniem powtórzeń jego wyrazów, nie tylko zbioru tych wyrazów. Na przykład ciąg stały $x_n=a$ ma punkt skupienia $a$, mimo że w zbiorze $\{a\}$ nie ma punktów różnych od $a$.

\begin{lemat}
\label{ma1:lem:subsequence-limit}
Jeżeli $(x_n)$ jest ciągiem w przestrzeni metrycznej $(X,d)$ zbieżnym do $a\in X$, to każdy jego podciąg jest zbieżny do $a$.
\end{lemat}
\begin{proof}
Dla ściśle rosnących dodatnich indeksów $n_k$ zachodzi $n_k\ge k$: dla $k=1$ mamy $n_1\ge1$, a z $n_k\ge k$ wynika $n_{k+1}\ge n_k+1\ge k+1$. Dla $\varepsilon>0$ wybieramy $N$ z $d(x_n,a)<\varepsilon$ dla wszystkich $n\ge N$. Gdy $k\ge N$, zachodzi $n_k\ge k\ge N$, zatem $d(x_{n_k},a)<\varepsilon$.
\end{proof}

\begin{twierdzenie}
\label{ma1:thm:cluster-criterion}
Niech $(x_n)$ będzie ciągiem w przestrzeni metrycznej $(X,d)$, a $a\in X$. Następujące warunki są równoważne:
\begin{enumerate}[label=(\roman*)]
\item $a$ jest punktem skupienia ciągu $(x_n)$;
\item dla każdego rzeczywistego $\varepsilon>0$ zbiór $\{n\in\N_+:d(x_n,a)<\varepsilon\}$ jest nieskończony;
\item dla każdego $k\in\N_+$ i $N\in\N_+$ istnieje $n\ge N$ taki, że $d(x_n,a)<1/k$.
\end{enumerate}
\end{twierdzenie}
\begin{proof}
Jeśli $x_{n_k}\to a$, to dla każdego $\varepsilon>0$ od pewnego indeksu $k$ wszystkie wyrazy podciągu należą do $B_d(a,\varepsilon)$. Indeksy $n_k$ są różne, więc odpowiedni zbiór indeksów jest nieskończony. Stąd (ii). Nieskończony podzbiór $\N_+$ nie jest ograniczony z góry przez żadną liczbę naturalną, ponieważ każdy początkowy odcinek $\{1,\ldots,N\}$ jest skończony. Wobec tego (ii) daje (iii) dla $\varepsilon=1/k$.
Załóżmy (iii). Wybieramy $n_1\ge1$ z $d(x_{n_1},a)<1$. Po wybraniu $n_k$ stosujemy (iii) dla $N=n_k+1$ i dokładności $1/(k+1)$, otrzymując $n_{k+1}>n_k$ z $d(x_{n_{k+1}},a)<1/(k+1)$. Możemy w każdym kroku wybierać najmniejszy dopuszczalny indeks. Dla $\varepsilon>0$ wybieramy $K$ z $1/K<\varepsilon$. Gdy $k\ge K$, mamy $d(x_{n_k},a)<1/k\le1/K<\varepsilon$. Otrzymaliśmy (i).
\end{proof}

\begin{twierdzenie}[Kryterium podciągów]
\label{ma1:thm:subsequence-criterion}
Niech $(x_n)$ będzie ciągiem w przestrzeni metrycznej $(X,d)$ i $a\in X$. Wówczas $x_n\to a$ wtedy i tylko wtedy, gdy każdy podciąg $(x_n)$ ma podciąg zbieżny do $a$.
\end{twierdzenie}
\begin{proof}
Konieczność wynika z lematu \ref{ma1:lem:subsequence-limit}. Załóżmy warunek dotyczący podciągów. Jeśli $x_n$ nie zbiega do $a$, to na mocy stwierdzenia \ref{ma1:prop:nonconvergence} istnieje $\varepsilon>0$ takie, że dla każdego $N$ można znaleźć $n\ge N$ z $d(x_n,a)\ge\varepsilon$. Wybierając kolejno najmniejsze takie indeksy, pierwszy co najmniej $1$, a następne większe od poprzednich, otrzymujemy podciąg $(x_{n_k})$ spełniający $d(x_{n_k},a)\ge\varepsilon$ dla każdego $k$. Każdy jego podciąg nadal spełnia tę nierówność, a zatem żaden nie zbiega do $a$. Sprzeczność.
\end{proof}

\begin{przyklad}
Ciąg $x_n=(-1)^n$ ma dwa punkty skupienia: $1$ i $-1$. Odpowiednie podciągi to $(x_{2k})$ i $(x_{2k-1})$. Innego punktu skupienia $a$ nie ma: jeśli $a\notin\{-1,1\}$, to $\varepsilon=\min(|a-1|,|a+1|)/2>0$ i żaden wyraz ciągu nie należy do $B(a,\varepsilon)$. Z twierdzenia \ref{ma1:thm:cluster-criterion} wynika teza. Ciąg $y_{2k}=0$, $y_{2k-1}=k$ ma tylko jeden punkt skupienia $0$, ale nie zbiega do niego, ponieważ dla dowolnego $N$ można wybrać nieparzysty indeks $2k-1\ge N$ z $|y_{2k-1}|\ge1$. Żaden punkt $a\ne0$ nie jest punktem skupienia $(y_n)$: dla $\varepsilon=\min(1,|a|/2)>0$ kula $B(a,\varepsilon)$ nie zawiera zera, a spośród dodatnich liczb całkowitych zawiera tylko skończenie wiele (każda jest mniejsza od $|a|+1$). Każda dodatnia liczba całkowita występuje w ciągu dokładnie raz, więc twierdzenie \ref{ma1:thm:cluster-criterion} wyklucza $a$ jako punkt skupienia. Jedyność punktu skupienia nie wystarcza więc do zbieżności.
\end{przyklad}

\section{Ciągi monotoniczne i twierdzenie Bolzana--Weierstrassa}
\begin{definicja}
\label{ma1:def:monotone-sequence}
Ciąg rzeczywisty $(a_n)$ nazywamy niemalejącym, jeśli $a_n\le a_{n+1}$ dla każdego $n\ge1$, a nierosnącym, jeśli $a_n\ge a_{n+1}$ dla każdego $n\ge1$. Ciąg monotoniczny to ciąg niemalejący lub nierosnący. Dla nierówności ścisłych używamy odpowiednio określeń rosnący i malejący.
\end{definicja}

\begin{twierdzenie}
\label{ma1:thm:monotone-convergence}
Niech $(a_n)$ będzie ciągiem rzeczywistym. Jeżeli jest niemalejący i ograniczony z góry, to zbiega do $\sup\{a_n:n\ge1\}$. Jeżeli jest nierosnący i ograniczony z dołu, to zbiega do $\inf\{a_n:n\ge1\}$.
\end{twierdzenie}
\begin{proof}
W pierwszym przypadku zbiór wyrazów jest niepusty i ograniczony z góry, więc z zupełności porządkowej $\R$ istnieje jego kres górny $s$. Dla $\varepsilon>0$ liczba $s-\varepsilon$ nie jest ograniczeniem górnym. Istnieje więc $N$ z $a_N>s-\varepsilon$. Z monotoniczności dla $n\ge N$ mamy $s-\varepsilon<a_N\le a_n\le s$, skąd $|a_n-s|<\varepsilon$. W drugim przypadku ciąg $(-a_n)$ jest niemalejący i ograniczony z góry, a $\sup\{-a_n:n\ge1\}=-\inf\{a_n:n\ge1\}$. Z pierwszej części i rachunku granic wynika teza.
\end{proof}
\begin{twierdzenie}
\label{ma1:thm:monotone-subsequence}
Każdy ciąg rzeczywisty ma podciąg monotoniczny.
\end{twierdzenie}
\begin{proof}
Niech $P=\{n\in\N_+:\forall m>n,\ a_m\le a_n\}$. Jeżeli $P$ jest nieskończony, wybieramy jego elementy w kolejności rosnącej $n_1<n_2<\cdots$; można to robić przez wybór najmniejszego elementu większego od poprzedniego. Definicja $P$ daje $a_{n_{k+1}}\le a_{n_k}$, więc podciąg jest nierosnący.
Jeśli $P$ jest skończony, wybieramy $n_1>\max P$, gdy $P\ne\varnothing$, a $n_1=1$, gdy $P=\varnothing$. Dla każdego $n_k\notin P$ istnieje $n_{k+1}>n_k$ z $a_{n_{k+1}}>a_{n_k}$, ponieważ jest to zaprzeczenie warunku definiującego $P$. Każdy nowy indeks również leży poza $P$. Otrzymujemy podciąg rosnący.
\end{proof}

\begin{twierdzenie}[Bolzana--Weierstrassa]
\label{ma1:thm:bolzano-weierstrass}
Każdy ograniczony ciąg w $\R^m$ dla $m\in\N_+$, wyposażonej w metrykę euklidesową, ma podciąg zbieżny. To samo zachodzi w $\C$ z metryką modułu.
\end{twierdzenie}
\begin{proof}
Dla $m=1$ wybieramy podciąg monotoniczny na mocy twierdzenia \ref{ma1:thm:monotone-subsequence}. Jest on ograniczony, gdyż każdy jego wyraz jest wyrazem ciągu wyjściowego. Twierdzenie \ref{ma1:thm:monotone-convergence} daje jego zbieżność.
Niech $m>1$ i załóżmy, że twierdzenie zachodzi w $\R^{m-1}$. Z nierówności $d_\infty\le d_2$ wynika ograniczoność pierwszych $m-1$ współrzędnych oraz ostatniej współrzędnej: jeśli $d_2(x^{(n)},p)<R$, to $|x_j^{(n)}-p_j|<R$ dla wszystkich $j,n$. Z założenia indukcyjnego istnieje podciąg $(x^{(n_k)})$, którego pierwsze $m-1$ współrzędnych zbiegają do pewnego punktu $\R^{m-1}$. Ostatnia współrzędna tego podciągu jest ograniczonym ciągiem rzeczywistym. Z przypadku jednowymiarowego wybieramy jego podciąg zbieżny, o indeksach $k_l$. W ciągu $(x^{(n_{k_l})})$ pierwsze $m-1$ współrzędnych nadal zbiegają do poprzednich granic, na mocy lematu \ref{ma1:lem:subsequence-limit}. Ostatnia współrzędna zbiega z konstrukcji. Twierdzenie \ref{ma1:thm:finite-coordinate-convergence} i porównanie metryk z przykładu \ref{ma1:ex:vector-metrics} dają zbieżność całego podciągu w $d_2$. Dla $\C$ stosujemy przypadek $m=2$, ponieważ $|z-w|$ jest odległością euklidesową odpowiednich par.
\end{proof}
\begin{przyklad}
Założenia o przestrzeni nie można zastąpić samą ograniczonością w dowolnej przestrzeni metrycznej. Niech $X=\N_+$ z metryką dyskretną. Ciąg $x_n=n$ jest ograniczony, ponieważ wszystkie odległości są nie większe od $1$. Każdy jego podciąg ma parami różne wyrazy, więc nie jest od pewnego indeksu stały. W metryce dyskretnej zbieżność wymaga stałości od pewnego indeksu. Żaden podciąg nie jest zatem zbieżny.
\end{przyklad}

\section{Ciągi Cauchy'ego}
\begin{definicja}
\label{ma1:def:cauchy-sequence}
Niech $(X,d)$ będzie przestrzenią metryczną. Ciąg $(x_n)_{n\ge1}$ nazywamy ciągiem Cauchy'ego, jeżeli
\[
 \forall\varepsilon\in\R_{>0}\quad\exists N\in\N_+\quad
 \forall m,n\ge N\quad d(x_m,x_n)<\varepsilon.
\]
\end{definicja}
Warunek dotyczy odległości wszystkich dostatecznie późnych par wyrazów, a nie tylko wyrazów kolejnych. Nie wymaga wskazania punktu granicznego.

\begin{twierdzenie}
\label{ma1:thm:cauchy-basic}
W dowolnej przestrzeni metrycznej zachodzą następujące własności:
\begin{enumerate}[label=(\roman*)]
\item każdy ciąg zbieżny jest ciągiem Cauchy'ego;
\item każdy ciąg Cauchy'ego jest ograniczony;
\item każdy podciąg ciągu Cauchy'ego jest ciągiem Cauchy'ego;
\item jeżeli ciąg Cauchy'ego ma podciąg zbieżny do $a$, to sam zbiega do $a$.
\end{enumerate}
\end{twierdzenie}
\begin{proof}
Dla (i), gdy $x_n\to a$, wybieramy $N$ z $d(x_n,a)<\varepsilon/2$ dla $n\ge N$. Wtedy dla $m,n\ge N$ zachodzi $d(x_m,x_n)\le d(x_m,a)+d(a,x_n)<\varepsilon$.
Dla (ii) wybieramy $N$ tak, aby $d(x_m,x_n)<1$ dla $m,n\ge N$. Ustalając $m=N$, otrzymujemy ograniczenie wszystkich późnych wyrazów. Dla
$R=1+\max(\{1\}\cup\{d(x_j,x_N):1\le j<N\})$ każdy wyraz należy do $B_d(x_N,R)$.
Dla (iii), jeśli $n_k$ są indeksami podciągu, to $n_k\ge k$ z dowodu lematu \ref{ma1:lem:subsequence-limit}. Gdy warunek Cauchy'ego zachodzi dla wszystkich indeksów co najmniej $N$, zachodzi również dla $n_k,n_l$ przy $k,l\ge N$.
Dla (iv) wybieramy $N$ tak, aby $d(x_m,x_n)<\varepsilon/2$ dla wszystkich $m,n\ge N$. Z $x_{n_k}\to a$ wybieramy $K$ z $d(x_{n_k},a)<\varepsilon/2$ dla $k\ge K$. Ustalmy $k\ge\max(K,N)$; wtedy $n_k\ge N$. Dla każdego $n\ge N$ mamy $d(x_n,a)\le d(x_n,x_{n_k})+d(x_{n_k},a)<\varepsilon$. Stąd $x_n\to a$.
\end{proof}
\begin{przyklad}
Warunek $d(x_{n+1},x_n)\to0$ nie wystarcza do warunku Cauchy'ego. W $\R$ połóżmy $s_n=\sum_{j=1}^n 1/j$. Wtedy $|s_{n+1}-s_n|=1/(n+1)\to0$. Tymczasem
\[
 s_{2n}-s_n=\sum_{j=n+1}^{2n}\frac1j
 \ge n\frac1{2n}=\frac12.
\]
Dla dowolnego $N$ wybieramy $n\ge N$, a indeksy $n,2n$ naruszają warunek Cauchy'ego dla $\varepsilon=1/2$.
\end{przyklad}

\section{Przestrzenie metryczne zupełne}
\begin{definicja}
\label{ma1:def:complete-metric}
Przestrzeń metryczną $(X,d)$ nazywamy zupełną, jeżeli każdy ciąg Cauchy'ego w $X$ jest zbieżny do pewnego punktu $X$.
\end{definicja}
Zupełność porządkowa z rozdziału \ref{ma1:ch:ordered-fields} i zupełność metryczna są różnymi pojęciami: pierwsze odnosi się do kresów, drugie do ciągów i odległości. Dla $\R$ wykażemy teraz zupełność metryczną, wykorzystując wcześniej udowodnioną zupełność porządkową.

\begin{twierdzenie}
\label{ma1:thm:real-metric-complete}
Przestrzeń $\R$ z metryką wartości bezwzględnej jest zupełna.
\end{twierdzenie}
\begin{proof}
Niech $(a_n)$ będzie ciągiem Cauchy'ego w $\R$. Z twierdzenia \ref{ma1:thm:cauchy-basic}(ii) jest ograniczony. Na mocy twierdzenia \ref{ma1:thm:bolzano-weierstrass} dla $m=1$ istnieje jego podciąg zbieżny do pewnego $a\in\R$. Twierdzenie \ref{ma1:thm:cauchy-basic}(iv) daje zbieżność całego ciągu do $a$. Każdy ciąg Cauchy'ego ma zatem granicę w $\R$.
\end{proof}
\begin{przyklad}
\label{ma1:ex:rational-incomplete}
Przestrzeń $\Q$ z metryką wartości bezwzględnej nie jest zupełna. Niech $r_n=\lfloor n\sqrt2\rfloor/n\in\Q$. Wtedy $r_n\le\sqrt2<r_n+1/n$, więc $|r_n-\sqrt2|<1/n$ i $r_n\to\sqrt2$ w $\R$. Z twierdzenia \ref{ma1:thm:cauchy-basic}(i) jest to ciąg Cauchy'ego w $\R$, a ponieważ odległości w $\Q$ są takie same, również w $\Q$. Jeśli miałby granicę $q\in\Q$, ta sama definicja dawałaby zbieżność do $q$ w $\R$. Jednoznaczność granicy implikowałaby $q=\sqrt2$, co przeczy niewymierności $\sqrt2$ z rozdziału \ref{ma1:ch:number-systems}. 
\end{przyklad}

\begin{twierdzenie}
\label{ma1:thm:complete-product}
Niech $m\in\N_+$, a $(X_j,d_j)$ dla $1\le j\le m$ będą niepustymi przestrzeniami metrycznymi. Produkt $X=\prod_{j=1}^m X_j$ z metryką $D_\infty=\max_j d_j$ jest zupełny wtedy i tylko wtedy, gdy każdy czynnik jest zupełny. Teza pozostaje prawdziwa dla $D_1=\sum_j d_j$.
\end{twierdzenie}
\begin{proof}
Jeśli każdy czynnik jest zupełny, a $(x^{(n)})$ jest ciągiem Cauchy'ego w $D_\infty$, to z $d_j(x_j^{(n)},x_j^{(k)})\le D_\infty(x^{(n)},x^{(k)})$ wynika warunek Cauchy'ego w każdym czynniku. Istnieją więc $a_j\in X_j$ z $x_j^{(n)}\to a_j$. Twierdzenie \ref{ma1:thm:finite-coordinate-convergence} daje zbieżność do $a=(a_1,\ldots,a_m)\in X$.
Odwrotnie, ustalmy $j$ i ciąg Cauchy'ego $(u_n)$ w $X_j$. Dla każdego $i\ne j$ wybierzmy stały punkt $p_i\in X_i$ (czynniki są niepuste). Niech $x^{(n)}$ ma $j$-tą współrzędną $u_n$, a pozostałe równe $p_i$. Wtedy $D_\infty(x^{(n)},x^{(k)})=d_j(u_n,u_k)$, więc jest to ciąg Cauchy'ego w $X$. Zupełność produktu daje granicę $a\in X$, a z nierówności $d_j(u_n,a_j)\le D_\infty(x^{(n)},a)$ wynika $u_n\to a_j$.
Dla $D_1$ korzystamy z $D_\infty\le D_1\le m D_\infty$. Jeśli ciąg jest Cauchy'ego w $D_\infty$, to wybór dokładności $\varepsilon/m$ daje warunek Cauchy'ego w $D_1$; odwrotna implikacja wynika z dolnego oszacowania. Obie metryki mają też te same ciągi zbieżne do danych punktów, na mocy twierdzenia \ref{ma1:thm:comparable-metrics}. Zupełność jest więc równoważna dla obu metryk.
\end{proof}

\begin{wniosek}
\label{ma1:cor:euclidean-complete}
Dla każdego $m\in\N_+$ przestrzeń $\R^m$ jest zupełna w metrykach $d_1,d_2,d_\infty$. Ciało $\C$ z metryką modułu również jest przestrzenią zupełną.
\end{wniosek}
\begin{proof}
Dla $d_1,d_\infty$ stosujemy twierdzenia \ref{ma1:thm:real-metric-complete} i \ref{ma1:thm:complete-product}. Z $d_\infty\le d_2\le\sqrt m\,d_\infty$ wynika, przez dokładnie ten sam argument z dokładnościami $\varepsilon/\sqrt m$ i $\varepsilon$, równoważność warunków Cauchy'ego dla $d_2$ i $d_\infty$. Zbieżność w obu metrykach jest równoważna z twierdzenia \ref{ma1:thm:comparable-metrics}. Dla $\C$ utożsamiamy moduł z $d_2$ dla $m=2$.
\end{proof}

\section{Zupełność podprzestrzeni}
\begin{definicja}
\label{ma1:def:closed-metric-subset}
Podzbiór $A$ przestrzeni metrycznej $(X,d)$ nazywamy domkniętym w $X$, jeżeli dla każdego $x\in X\setminus A$ istnieje $r>0$ w $\R$ takie, że $B_d(x,r)\cap A=\varnothing$.
\end{definicja}
\begin{lemat}[Ciągowe kryterium domkniętości]
\label{ma1:lem:closed-sequential}
Podzbiór $A$ przestrzeni metrycznej $(X,d)$ jest domknięty wtedy i tylko wtedy, gdy dla każdego ciągu $(a_n)$ w $A$ i każdego $x\in X$ z $a_n\to x$ w $X$ wynika $x\in A$.
\end{lemat}
\begin{proof}
Jeśli $A$ jest domknięty i $a_n\to x$ dla ciągu w $A$, załóżmy $x\notin A$. Istnieje wówczas $r>0$ z $B_d(x,r)\cap A=\varnothing$. Z definicji granicy od pewnego indeksu $a_n\in B_d(x,r)$, wbrew $a_n\in A$.
Odwrotnie, jeśli $A$ nie jest domknięty, istnieje $x\notin A$ takie, że dla każdego $r>0$ kula $B_d(x,r)$ przecina $A$. Dla każdego $n\ge1$ wybieramy $a_n\in A\cap B_d(x,1/n)$ (korzystamy z wyboru przeliczalnego). Dla $\varepsilon>0$ wybieramy $N$ z $1/N<\varepsilon$. Dla $n\ge N$ zachodzi $d(a_n,x)<1/n\le1/N<\varepsilon$, więc $a_n\to x$, lecz $x\notin A$. To przeczy warunkowi ciągowemu.
\end{proof}

\begin{twierdzenie}
\label{ma1:thm:complete-closed}
Niech $(X,d)$ będzie przestrzenią metryczną, a $A\subseteq X$ z metryką indukowaną. Jeżeli $A$ jest zupełna, to jest domknięta w $X$. Jeśli $X$ jest zupełna, to każdy domknięty podzbiór $A$ jest zupełny.
\end{twierdzenie}
\begin{proof}
Załóżmy zupełność $A$ i niech $(a_n)$ będzie ciągiem w $A$ zbieżnym w $X$ do $x\in X$. Jest to ciąg Cauchy'ego w $X$ na mocy twierdzenia \ref{ma1:thm:cauchy-basic}(i), a odległości w $A$ są takie same. Z zupełności $A$ istnieje $a\in A$ z $a_n\to a$ w $A$, a zatem również w $X$. Jednoznaczność granicy daje $x=a\in A$. Lemat \ref{ma1:lem:closed-sequential} dowodzi domkniętości.
Jeśli $X$ jest zupełna i $A$ domknięta, niech $(a_n)$ będzie ciągiem Cauchy'ego w $A$. Jest on ciągiem Cauchy'ego w $X$, więc ma granicę $x\in X$. Z domkniętości i lematu \ref{ma1:lem:closed-sequential} wynika $x\in A$. Ponieważ metryka jest ograniczeniem $d$, zbieżność zachodzi też w $A$. Dla zbioru pustego oba warunki są spełnione w sposób pusty: nie istnieją w nim ciągi.
\end{proof}
\begin{przyklad}
Dla $a,b\in\R$ z $a\le b$ przedział $[a,b]$ jest domknięty w $\R$: jeśli $x<a$, promień $(a-x)/2$ daje kulę rozłączną z $[a,b]$; dla $x>b$ wybieramy $(x-b)/2$. Zatem $[a,b]$ jest zupełny. Natomiast $(0,1]$ nie jest zupełny, bo ciąg $1/n$ jest ciągiem Cauchy'ego w tej podprzestrzeni, lecz nie ma w niej granicy. Wynika to z twierdzenia \ref{ma1:thm:cauchy-basic}(i) zastosowanego w $\R$ i jednoznaczności granicy.
\end{przyklad}

\section{Zupełność przestrzeni funkcji ograniczonych}
\begin{twierdzenie}
\label{ma1:thm:bounded-functions-complete}
Niech $S$ będzie niepustym zbiorem, a $K=\R$ lub $K=\C$. Niech $\mathcal B(S,K)$ oznacza zbiór funkcji $f:S\to K$ o ograniczonym zbiorze wartości. Przestrzeń $\mathcal B(S,K)$ z metryką
$d_\infty(f,g)=\sup_{t\in S}|f(t)-g(t)|$ jest zupełna.
\end{twierdzenie}
\begin{proof}
Dla $f,g\in\mathcal B(S,K)$ istnieją rzeczywiste $M,L>0$ z $|f(t)|\le M$ i $|g(t)|\le L$ dla wszystkich $t$. Zatem $|f(t)-g(t)|\le M+L$ i supremum istnieje. Nieujemność i symetria są własnościami modułu; zerowanie supremum wymusza $|f(t)-g(t)|=0$ dla wszystkich $t$, a dla $f=g$ supremum wynosi zero. Nierówność trójkąta wynika z nierówności $|f(t)-h(t)|\le d_\infty(f,g)+d_\infty(g,h)$ dla każdego $t$ i przejścia do supremum.

Niech $(f_n)$ będzie ciągiem Cauchy'ego w tej metryce. Dla ustalonego $t\in S$ z $|f_m(t)-f_n(t)|\le d_\infty(f_m,f_n)$ wynika, że $(f_n(t))$ jest ciągiem Cauchy'ego w $K$. Z zupełności $\R$ lub $\C$ istnieje dokładnie jedna granica $f(t)\in K$, więc definiujemy funkcję $f:S\to K$.

Dla $\varepsilon>0$ wybieramy $N$ takie, że $d_\infty(f_m,f_n)<\varepsilon/2$ dla $m,n\ge N$. Ustalmy $n\ge N$ oraz $t\in S$. Dla każdego $m\ge N$ zachodzi $|f_m(t)-f_n(t)|<\varepsilon/2$. Z twierdzenia \ref{ma1:thm:distance-limit} (drugi ciąg stały) odległości te zbiegają do $|f(t)-f_n(t)|$. Zachowanie nierówności nieostrych w granicy (twierdzenie \ref{ma1:thm:sequence-squeeze}(i)) daje $|f(t)-f_n(t)|\le\varepsilon/2$. Oszacowanie obowiązuje dla wszystkich $t$ i wszystkich $n\ge N$.

Dla $\varepsilon=2$ wybieramy taki indeks $N$ i ustalamy $n=N$. Ponieważ $f_N$ jest ograniczona, istnieje $M>0$ takie, że $|f_N(t)|\le M$ dla wszystkich $t$. Otrzymujemy $|f(t)|\le M+1$, czyli $f\in\mathcal B(S,K)$. Dla dowolnego $\varepsilon>0$ wcześniejsze oszacowanie daje teraz $d_\infty(f,f_n)\le\varepsilon/2<\varepsilon$ dla wszystkich $n\ge N$. Zatem $f_n\to f$ w metryce supremum.
\end{proof}
Przestrzeń ograniczonych ciągów rzeczywistych z rozdziału \ref{ma1:ch:metric-sequences} jest szczególnym przypadkiem dla $S=\N_+$. Jest zupełna, mimo że ograniczony ciąg jej elementów nie musi mieć zbieżnego podciągu. Własności te są różne.

\begin{uwaga}
W rozdziale wykorzystywaliśmy zupełność $\R$ do istnienia granic i kresów. W następnym rozdziale opiszemy dolną i górną granicę ciągu rzeczywistego za pomocą kresów zbiorów jego dalszych wyrazów. Nie utożsamiamy ich z granicą zwykłą bez odpowiednich założeń.
\end{uwaga}
